\section{JUSTIFICACIÓN}

\subsection{Justificación Social y Ciudadana}
El acceso a información transparente y oportuna respecto al costo de bienes de primera necesidad y medicamentos esenciales constituye un derecho fundamental para la población. En coyunturas económicas caracterizadas por presiones inflacionarias o variaciones en la disponibilidad de insumos importados, la ciudadanía boliviana suele experimentar incertidumbre al adquirir productos para la salud o la nutrición básica. 

La implementación de este sistema democratiza el acceso a la información comercial al consolidar en un punto único las ofertas de cadenas minoristas y farmacéuticas tradicionalmente dispersas. Asimismo, la incorporación de un módulo de georreferenciación que ubica los 83 centros de salud públicos y de atención primaria del departamento de La Paz, calculando distancias geodésicas inmediatas desde la posición del usuario, aporta un valor social tangible al orientar a la población hacia el establecimiento más cercano en situaciones de requerimiento médico.

\subsection{Justificación Económica y de Transparencia}
Desde una perspectiva microeconómica, los mercados minoristas funcionan de manera más eficiente cuando se reduce la asimetría de información entre compradores y oferentes. La capacidad de contrastar el precio de un mismo fármaco o alimento entre Hipermaxi, Farmacorp y Farmacias Chávez en cuestión de segundos empodera a las familias bolivianas para optimizar su presupuesto mensual y generar ahorros directos que, según las mediciones del sistema, oscilan entre el 15\% y el 45\% en productos seleccionados. 

En el ámbito de las compras del Estado, el análisis cuantitativo de las contrataciones públicas registradas en el SICOES contribuye a la transparencia gubernamental, permitiendo examinar de manera sistemática los volúmenes presupuestarios comprometidos en licitaciones públicas de La Paz.

\subsection{Justificación Técnica y de Ingeniería}
El desarrollo del proyecto representa una contribución relevante a las prácticas modernas de ingeniería de datos y minería web en entornos de alta hostilidad WAF. Frente al agotamiento de las técnicas tradicionales de scraping mediante scripts estáticos, la metodología propuesta articula la instrumentación del protocolo CDP sobre navegadores Chromium persistentes, la resolución programada de retos anti-bot, el particionamiento físico de series temporales en bases de datos relacionales serverless y el almacenamiento inmutable de objetos multimedia en la nube. 

La integración pionera de modelos densos de embeddings de 1024 dimensiones operando junto a índices vectoriales HNSW en PostgreSQL demuestra la viabilidad técnica de implementar arquitecturas híbridas de recuperación de información que combinan el rigor relacional SQL con la flexibilidad de la inteligencia artificial moderna.

\subsection{Justificación Académica e Institucional}
En el marco académico de la Carrera de Informática de la Universidad Mayor de San Andrés (UMSA), el proyecto responde a las exigencias formativas de la asignatura de Minería de Datos. El trabajo consolida en una única solución práctica y formal los tres pilares conceptuales de la Minería Web (\textit{Content}, \textit{Structure} y \textit{Usage Mining}), proveyendo un marco metodológico reproducible que servirá como base de consulta y extensión para futuras investigaciones en ciencia de datos, sistemas distribuidos y analítica de datos en Bolivia.
